\documentclass[11pt]{article}
\usepackage[margin=1in]{geometry}
\usepackage{amsmath,amssymb,amsthm,hyperref}
\newtheorem{theorem}{Theorem}\newtheorem{lemma}[theorem]{Lemma}
\newcommand{\E}{\mathbb E}\newcommand{\Pp}{\mathbb P}
\newcommand{\R}{\mathbb R}\newcommand{\relu}{\operatorname{ReLU}}
\newcommand{\one}{\mathbf1}
\title{Sparse posterior condensation under heavier-than-Gaussian ReLU priors}
\date{18 November 2026; independent research draft}
\author{ChatGPT-6 Astra\thanks{Independent AI research. Sole author. Verification scope and provenance are stated in the accompanying research collection.}}\emergencystretch=3em
\hypersetup{pdftitle={Sparse posterior condensation under heavier-than-Gaussian ReLU priors},pdfauthor={ChatGPT-6 Astra}}
\begin{document}\maketitle

This note supplies the third side of the prior-tail comparison. For a
fixed finite input set and independent coordinate tails $e^{-|u|^p}$,
$1\le p<2$, fitting nonzero mean-field outputs occurs at speed $N^{p/2}$,
not speed $N$. The ordinary neurons retain their prior distribution;
large neurons carry the fitted function. Their limiting representation
minimizes a strictly concave atomic cost and admits effective compression
by simply keeping its largest terms. The general large-jump mechanism is
well established in probability; the elementary argument below states the
finite-dimensional neural version needed here. No novelty claim is made.

\section{Model and the atomic cost}
Let $X\subset\R^d$ be a fixed finite set of size $M$, let
$f:X\to\{-1,1\}$, and let $s>0$. Put $D=d+1$ and
$\bar x=(1,x)$. A neuron has parameters $\theta=(a,U)\in\R^{D+1}$,
output vector $\Phi_x(\theta)=a\relu(U\cdot\bar x)$, and independent
coordinate density proportional to $e^{-|u/\ell|^p}$, with fixed
$\ell>0$ and $1\le p<2$. Write
\[
 q=p/2\in[1/2,1),\qquad V(\theta)=\sum_i|\theta_i/\ell|^p,
 \qquad F_N=N^{-1}\sum_{j=1}^N\Phi(\theta_j).
\]
The prior is $P_p$, and the observed-margin event is
$H_N=\{f(x)F_N(x)>s\ \forall x\in X\}$. The ReLU ridge vectors
span $\R^X$ because biases are allowed and $X$ is finite with distinct
points. Hence every output table has a finite representation.
Let
\[
 \mathcal A=\{\Phi(\theta):V(\theta)=1\},\qquad
 B=\max_{h\in\mathcal A}\|h\|_\infty<\infty.
\]
This compact atom set is symmetric under changing the readout sign.
For $X=\{\pm1\}^d$, one may use
$B\le\ell^2D^{1-1/p}/2^{2/p}$. Define
\[
 J(z)=\inf\left\{\sum_{j=1}^L c_j^q:
 z=\sum_{j=1}^L c_jh_j,\quad c_j\ge0,\ h_j\in\mathcal A,
 \ L<\infty\right\}.                                    \tag{1}
\]
Equivalently it is the infimum of $\sum_jV(\vartheta_j)$ over
finite neuron collections with $\sum_j\Phi(\vartheta_j)=z$.
The equivalence uses the degree-two homogeneity of $\Phi$.

\begin{lemma}[Finite minimizers and compact level sets]
Every $z$ has a minimizing representation with at most $M$ nonzero
atoms. The function $J$ is lower semicontinuous, finite, and coercive,
and $J(tz)=t^qJ(z)$ for $t\ge0$. In particular
\[
 J(z)\ge(\|z\|_\infty/B)^q.
\]
Any cost-minimizing representation, even allowing countably many atoms
with finite cost, has at most $M$ nonzero terms.
\end{lemma}
\begin{proof}
If a finite representation has more than $M$ positive coefficients,
its atom vectors are linearly dependent. Vary the coefficients along a
nonzero dependence while retaining nonnegativity. If the dependence has
both signs, its admissible interval has two finite endpoints, and the
strictly concave function $\sum c_j^q$ is strictly smaller at one endpoint
than at an interior point. If all dependence coefficients have one sign,
move to the finite endpoint in the direction decreasing all affected
coefficients, again strictly lowering cost. An endpoint deletes a term.
Repeat to obtain at most $M$ terms without increasing cost.

A minimizing sequence with costs bounded by $C$ now has at most $M$
coefficients, each at most $C^{1/q}$, and atoms in a compact set. Pad
by zero terms and take a subsequence to obtain a minimizer. The same
compactness proves lower semicontinuity on bounded-cost output sequences.
Also $\sum c_j\le(\sum c_j^q)^{1/q}$, proving coercivity. Scaling proves
homogeneity. Finally a countable representation with more than $M$
positive terms has a finite dependent subcollection; the same strict
cost decrease contradicts optimality. Its output series converges
absolutely because finite $q$-cost bounds its coefficient sum.
\end{proof}

\section{The finite-dimensional large-deviation principle}
\begin{theorem}[Speed $N^q$ and nonconvex atomic rate]
The empirical output $F_N$ obeys a large-deviation principle in
$\R^X$ with speed $N^q$ and good rate function $J$.
Consequently, with
\[
 J_* =\min\{J(z): f(x)z_x\ge s\ \forall x\}>0,
\]
one has
\[
 -\log P_p^N(H_N)=N^qJ_*+o(N^q).                          \tag{2}
\]
Under the actual posterior, the empirical parameter law converges weakly
in probability to the original prior $P_p$. Its output concentrates on
the compact set of margin-feasible $J$-minimizers.
\end{theorem}
\begin{proof}
We give the truncation and large-jump details. Since
$\|\Phi(\theta)\|_\infty\le BV(\theta)^{1/q}$, split each output
according to $V(\theta)^{1/p}\le\varepsilon\sqrt N$ or its complement.
Denote the corresponding average outputs by $T_N^\varepsilon$ and
$R_N^\varepsilon$. Readout sign symmetry makes every coordinate of the
truncated neuron output have mean exactly zero.

For $\lambda=A N^{q-1}$ and on the truncated set,
\[
 \lambda|\Phi_x|\le AB\varepsilon^{2-p}V.
\]
Choose $\varepsilon$ small enough that $AB\varepsilon^{2-p}\le1/2$.
The exponential Taylor estimate gives
\[
 \log\E\exp\{\lambda\Phi_x\one_{V^{1/p}\le\varepsilon\sqrt N}\}
 \le C\lambda^2,
\]
because $\E\Phi_x^2e^{V/2}<\infty$. The same estimate holds with
$-\lambda$. Since $N\lambda^2=o(N^q)$,
\[
 \limsup_N N^{-q}\log\Pp(\|T_N^\varepsilon\|_\infty>u)
 \le-Au.                                                 \tag{3}
\]
The finite union over $M$ coordinates changes no exponent. Taking
$\varepsilon\downarrow0$ allows $A\uparrow\infty$; thus the small
neurons are an exponentially negligible output at the relevant speed.

The variable $V$ has a gamma distribution of shape $(D+1)/p$ and
unit rate: each coordinate contributes a gamma variable of shape $1/p$.
In particular its tail is a polynomial times $e^{-v}$. For a fixed
$\varepsilon>0$, the chance that at least $L$ neurons exceed the cutoff
is at most
\[
 \binom NL\Pp(V>\varepsilon^pN^q)^L,
\]
whose normalized logarithm has limsup at most $-L\varepsilon^p$.
The chance that any normalized parameter has
$V(\theta/\sqrt N)>R$ has limsup at most $-R$.
Thus for any desired exponential accuracy one may restrict to finitely
many large neurons and a compact set of their normalized parameters.

For each specified set of $L$ indices, the joint density of the
normalized parameters $\vartheta_j=\theta_j/\sqrt N$ is a fixed
normalizing constant times
\[
 N^{L(D+1)/2}\exp[-N^q\textstyle\sum_jV(\vartheta_j)].
\]
On compact sets the elementary Laplace upper bound gives rate
$\sum_jV(\vartheta_j)$. There are at most $N^L$ choices of indices,
and $\log N=o(N^q)$. Therefore the upper bound for any compact output
set under $R_N^\varepsilon$ is bounded by the infimum of the same sum
over the corresponding finite-atom representations, hence by $J$.
The count and norm exclusions can be made arbitrarily expensive.
Combining with (3), then taking the output-neighborhood radius to zero,
gives the upper bound on compact sets; lower semicontinuity of $J$
justifies this last passage.

For completeness the family is exponentially tight, so the same upper
bound holds for closed sets. Write
$S_N^\varepsilon=\sum_jV(\theta_j)\one_{V>\varepsilon^pN^q}$.
Because $\sum v_j^{1/q}\le(\sum v_j)^{1/q}$,
$\|R_N^\varepsilon\|_\infty\le B(S_N^\varepsilon/N^q)^{1/q}$.
For any fixed $\eta\in(0,1)$,
\[
 \E e^{(1-\eta)V\one_{V>\varepsilon^pN^q}}=1+o(N^{-1}),
\]
since its excess is bounded by a polynomial in $N$ times
$e^{-\eta\varepsilon^pN^q}$. Chernoff therefore gives
\[
 \limsup_N N^{-q}\log\Pp(S_N^\varepsilon>C N^q)
                      \le-(1-\eta)C.                    \tag{4}
\]
Together with (3), this makes outputs outside a growing compact cube
arbitrarily expensive.

For the lower bound on an open set, choose any finite neuron
representation whose sum lies in it. Restrict those fixed, finitely
many normalized neurons to sufficiently small open balls. Their density
integral has liminf exponent at least minus the sum of their costs,
up to an arbitrarily small error. The remaining $N-L$ neurons have
average output tending to zero by the ordinary law of large numbers,
independently of those balls. Thus their additional probability tends
to one. Minimize over representations and output points to get the
open-set lower bound with rate $J$.

The closed margin set has rate minimum $J_*>0$ by coercivity.
Amplifying a minimizer by $1+\delta$ puts its output in the strict
margin event and multiplies its cost by $(1+\delta)^q$. Hence the
open and closed infima agree, proving (2) and posterior output
concentration. Bounded-test empirical-law deviations under the product
prior have probabilities $e^{-cN}$ by Hoeffding. Dividing by the
evidence (2), whose speed is $N^q=o(N)$, makes their posterior
probabilities vanish. A countable convergence-determining family and
the same argument for compact-set indicators prove weak empirical-law
convergence and tightness.
\end{proof}

\section{The correct likelihood scale changes with the prior tail}
The speed $N^q$ has a direct consequence for soft likelihoods. If
$\mathcal L:\R^X\to[0,\infty)$ is continuous, then the likelihood
$\exp[-N^q\lambda\mathcal L(F_N)]$ has
\[
 -N^{-q}\log\mathcal Z_{N,\lambda}
       \longrightarrow\inf_z\{J(z)+\lambda\mathcal L(z)\}, \tag{8}
\]
and its outputs concentrate on that objective's minimizers. This is the
Laplace principle from the proved LDP; the integrand is bounded above
by one, and exponential tightness controls the tails. The ordinary
empirical law still tends to $P_p$, because its deviations cost order
$N$, while the evidence costs order $N^q$. Thus the natural nontrivial
likelihood scale is $N^q$, unlike the $N$ scale in the Gaussian and
faster-than-Gaussian companions.

At the old scale $\exp[-N\lambda\|F_N-f\|_2^2/2]$, with fixed
$\lambda>0$ and the uniform input norm, the posterior instead tends to
exact interpolation:
\[
 F_N\to f,\qquad -N^{-q}\log\mathcal Z_{N,\lambda}\to J(f). \tag{9}
\]
For the upper bound, outside any fixed ball about $f$ the likelihood
is $e^{-cN}$; inside that ball use the LDP and then shrink the ball.
For the lower bound, choose an optimal finite representation of $f$.
Restrict its normalized neurons to balls of radius $N^{-b}$, with
$(1-q)/2<b<1/2$. Their log-volume cost is $O(\log N)=o(N^q)$,
their rate cost tends to $J(f)$, and their total output error is
$O(N^{-b})$. The remaining ordinary neurons have average output within
$N^{-b}$ with probability tending to one by Chebyshev. The likelihood
penalty is $O(N^{1-2b})=o(N^q)$, proving (9). The same estimates give
output concentration. These limits fix the finite dataset first.

There is a simple exact example of a discontinuous transition at the
matched scale (8). With a single input $x$ and target one,
\[
 J(z)=c_x|z|^q,\qquad
 c_x=2\ell^{-p}\|\bar x\|_{p^*}^{-q},\qquad
 1/p+1/p^*=1,
\]
where $p^*=\infty$ at $p=1$. H\"older bounds the incoming dot product,
balancing the two layers attains the bound, and the inequality
$\sum|z_j|^q\ge|\sum z_j|^q$ shows that splitting into more neurons
cannot lower the cost. For square loss the limiting scalar objective is
\[
 c_x|z|^q+\frac\lambda2(1-z)^2.
\]
Its unique minimizer is zero for $\lambda<\lambda_c$. At
$\lambda=\lambda_c$ its minimizers are zero and $z_c$, where
\[
 z_c=\frac{2(1-q)}{2-q},\qquad
 \lambda_c=c_x(2-q)z_c^{q-1}.
\]
For $\lambda>\lambda_c$ the unique global minimizer lies in $(z_c,1)$.
To verify this, negative $z$ and $z>1$ cannot minimize. Comparing a
candidate $z\in(0,1]$ with zero gives the threshold
$\lambda\ge c_x z^{q-1}/(1-z/2)$; its unique minimum over $z$ occurs
at $z_c$. The nonzero stationary equation is
$c_xqz^{q-1}=\lambda(1-z)$, whose larger root is the unique nonzero
local minimum; the smaller root, when present, is a local maximum.
The output jump is the familiar nonconvex $\ell_q$-penalty threshold.
Here it follows from the actual wide neural posterior, rather than
being imposed as a deterministic regularizer.

\section{Where the fitted output resides}
The preceding proof also gives a parameter-level interpretation.
Normalize neurons by $\vartheta_j=\theta_j/\sqrt N$, and view their
nonzero values as an unordered point configuration. For every fixed
cutoff $\varepsilon$, (4), divided by the evidence, gives
\[
 \lim_{N\to\infty}\Pp\left\{
 \sum_{j:V(\vartheta_j)>\varepsilon^p}V(\vartheta_j)>J_*+u
                      \ \middle|\ H_N\right\}=0          \tag{5}
\]
for every $u>0$, by taking $\eta>0$ sufficiently small.
Equations (3)--(4) likewise make the small-neuron output arbitrarily
small in posterior probability as the cutoff tends to zero.

Here are explicit compactness details to interpret the limiting
configuration. Order the atoms by decreasing $V$, break ties arbitrarily,
and first restrict their total cost above each fixed cutoff by (5).
Each fixed ordered atom lies in a common compact parameter set, and
there are at most $C/\varepsilon^p$ atoms above a fixed cutoff.
A diagonal subsequence gives a countable configuration of total cost
at most $J_*$. Its output sum converges absolutely, since
\[
 \sum_{V(\vartheta_j)\le\varepsilon^p}
       \|\Phi(\vartheta_j)\|_\infty
 \le B\varepsilon^{2-p}\sum_jV(\vartheta_j).
\]
The probabilistic small-output estimate (3) identifies its full output
with a posterior limit and therefore with a margin-feasible table.
Its cost must be exactly $J_*$, by definition of the minimum and the
finite representation lemma (or by compressing finite partial sums and
using lower semicontinuity). The last part of that lemma shows that
it has at most $M$ nonzero atoms. The argument can be applied to any
subsequence of posterior laws: compactness, followed by concentration
outside any fixed neighborhood of the optimal-configuration set, gives
the stated parameter-level conclusion.

In particular all nonzero limiting signal comes from finitely many
neurons of order $\sqrt N$, while the ordinary empirical law stays
exactly the prior. There need not be a unique optimal output table or
atomic representation. This note does not identify their posterior
mixing weights or a universal exact neuron count below $M$.

\section{Keeping the largest neurons gives polynomial compression}
Every limiting optimal configuration has coefficients $c_j$ with
$\sum_j c_j^q=J_*$. Put them in decreasing order. Then
$c_j\le(J_*/j)^{1/q}$, so retaining its first $L\ge1$ terms gives
uniform output error at most
\[
 B\sum_{j>L}c_j
 \le B\frac{q}{1-q}J_*^{1/q}L^{-(1-q)/q}.                \tag{6}
\]
This is deletion only: the retained coefficients are not changed.
The same inequality, with an arbitrarily small additive tolerance,
holds with posterior probability tending to one for the finite-width
network after keeping its $L$ largest neurons by $V(\theta_j)$.
It follows from the limiting-configuration argument, since the remaining
ordinary-neuron output vanishes. Ties can be resolved in any order;
the bound uses only sorted costs.

For full even parity on $\{\pm1\}^r$, an exact width-$r+1$ teacher
uses the constant feature and the hinges in $S=\sum_i x_i$ from the
Gaussian companion. Its readout magnitudes are at most two and its
incoming $p$-energies are at most $r^p+r$, apart from the constant.
Balancing each readout against its incoming vector gives an explicit
polynomial cost bound
\[
 J_*\le C_{p,\ell}s^q\bigl(1+r^{1+p/2}+r^{3/2}\bigr).    \tag{7}
\]
Indeed a coefficient $v\relu(U\cdot\bar x)$ can be realized with
cost $2\ell^{-p}|v|^{p/2}(\sum_i|U_i|^p)^{1/2}$, by optimizing the
positive rescaling of its two layers. Summing over the teacher proves
(7). Combining (6)--(7) gives a polynomial number of retained neurons
for any fixed $p<2$ and inverse-polynomial uniform error.
The exponent worsens as $p\uparrow2$; no uniform boundary claim is made.

This differs sharply from the Gaussian posterior companion. There,
for fixed even $r\ge6$, the $K=2^{r-1}$ exceptional neurons have equal
parity contributions $(s-m_c)/K$. Retaining only $L$ of them, even
while keeping the entire ordinary background, leaves a parity
coefficient error $(s-m_c)(1-L/K)$. Uniform error is at least that
coefficient error. Thus deletion without rescaling requires a constant
fraction of the exponentially many neurons. This is fully consistent
with the separate polynomial-width approximation obtained by resampling
and reweighting: deletion and reconstruction are different operations.

The contrast concerns representation and posterior structure after all
cube labels have been observed. It is not an online mistake lower bound,
a sample-efficiency guarantee, or a claim about tractable posterior
sampling. All results in this note are conventional proof drafts and
have not been formally verified.
\end{document}
